\notechapter{4}{Mathematical structure and proposed bulk-silicon comparison of tight-binding reduction routes}

\textbf{Status: mathematical framework and proposed numerical experiment.} No
bulk-silicon production calculation, accepted Wannier operator, fitted
tight-binding model, convergence result, or route comparison is claimed in
this note. The physical and numerical parent remain governed by the applicable
specifications and by Chapters~\ref{ch:physical-problem} and
\ref{ch:first-principles-parent}. The representation, alignment, and reduction
requirements remain governed by Chapters~\ref{ch:representations-and-alignment}
and~\ref{ch:model-reduction}.

\section*{Purpose}

The bulk-silicon program contains two proposed routes from a retained
Kohn--Sham parent to an orthogonal $sp^3s^{\ast}$ tight-binding model:
\begin{align}
    \text{direct route:}\qquad
    &\mathrm{KS\text{-}DFT}
    \longrightarrow
    \mathrm{TB},
    \\
    \text{operator-mediated route:}\qquad
    &\mathrm{KS\text{-}DFT}
    \longrightarrow
    \mathrm{Wannier}
    \longrightarrow
    \mathrm{TB}.
\end{align}
The routes share a parent but do not perform the same mathematical operations.
The direct route is initially a constrained spectral fit. The mediated route
constructs a retained localized representation and then compares aligned
operators. Their outputs cannot be compared responsibly until the state
spaces, maps, gauges, energy references, model classes, and fitting objectives
have been declared.

This note first states the mathematical issues that determine which
comparisons are meaningful. It then describes a proposed numerical experiment
for the bulk-silicon pilot. The experiment is intended to compare
KS-DFT-to-TB, KS-DFT-to-Wannier, and Wannier-to-TB relations without treating
them as three instances of one undifferentiated matrix norm.

\section*{Part I: Mathematical issues}

\subsection*{Parent Bloch fibers and retained subspaces}

Let $\hat H_{\mathrm{KS}}(\mathbf k)$ denote a Kohn--Sham Bloch-fiber
operator at crystal momentum $\mathbf k$. Let
$\mathcal S(\mathbf k)$ be an $M$-dimensional retained subspace with
orthogonal projector $\hat P(\mathbf k)$. The retained parent operator is
\begin{equation}
    \hat H_{\mathrm{KS}}^{(P)}(\mathbf k)
    =
    \left.
        \hat P(\mathbf k)
        \hat H_{\mathrm{KS}}(\mathbf k)
        \hat P(\mathbf k)
    \right|_{\mathcal S(\mathbf k)}.
\end{equation}
The exact band count, reciprocal-space domain, projection choices, and energy
windows are not fixed by this note. They must be inherited from the accepted
target-subspace specification.

A list of selected eigenvalues does not by itself define
$\mathcal S(\mathbf k)$. Degeneracies, crossings, and entanglement make the
projector the stable mathematical object. If disentanglement is used, it
selects a subspace from a larger candidate space and can therefore introduce a
subspace-selection discrepancy. A later unitary gauge transformation within a
fixed retained subspace does not change that projector.

The parent differential operator, the retained operator, and any finite matrix
representation remain distinct:
\begin{equation}
    \hat H_{\mathrm{KS}}(\mathbf k),
    \qquad
    \hat H_{\mathrm{KS}}^{(P)}(\mathbf k),
    \qquad
    \mathbf H_{\mathrm{KS}}^{(P)}(\mathbf k).
\end{equation}
Only the last object is directly available to a Frobenius-norm calculation,
and only after its basis and normalization have been declared.

\subsection*{From the retained KS operator to a Wannier representation}

Let
\begin{equation}
    \mathbf J_{\mathrm W}(\mathbf k):
    \mathbb C^M
    \longrightarrow
    \mathcal S(\mathbf k)
\end{equation}
be an isometric frame map for the selected Wannier gauge, so that
\begin{equation}
    \mathbf J_{\mathrm W}(\mathbf k)^\dagger
    \mathbf J_{\mathrm W}(\mathbf k)
    =
    \mathbf I_M.
\end{equation}
The retained Hamiltonian in Wannier-gauge Bloch coordinates is
\begin{equation}
    \mathbf H_{\mathrm W}(\mathbf k)
    =
    \mathbf J_{\mathrm W}(\mathbf k)^\dagger
    \hat H_{\mathrm{KS}}^{(P)}(\mathbf k)
    \mathbf J_{\mathrm W}(\mathbf k).
\end{equation}
For a fixed retained subspace and a complete representation, this is a change
of coordinates rather than a tight-binding approximation. Localization chooses
a gauge within the subspace. It does not by itself discard rank or hopping
blocks.

On a uniform grid of $N_{\mathbf k}$ points, the corresponding real-space
blocks use the declared Fourier convention
\begin{equation}
    \mathbf H_{\mathrm W}(\mathbf R)
    =
    \frac{1}{N_{\mathbf k}}
    \sum_{\mathbf k}
    e^{-i\mathbf k\cdot\mathbf R}
    \mathbf H_{\mathrm W}(\mathbf k).
\end{equation}
The inverse transform is
\begin{equation}
    \mathbf H_{\mathrm W}(\mathbf k)
    =
    \sum_{\mathbf R}
    e^{i\mathbf k\cdot\mathbf R}
    \mathbf H_{\mathrm W}(\mathbf R)
\end{equation}
when all represented blocks and matching conventions are retained. Discarding
blocks beyond a chosen range is a separate truncation.

The KS-DFT-to-Wannier relation must therefore distinguish at least:
\begin{enumerate}
    \item retained-subspace selection or disentanglement;
    \item gauge selection and localization;
    \item finite reciprocal-grid interpolation;
    \item any omitted real-space blocks; and
    \item the numerical fidelity of the producing calculation.
\end{enumerate}
A band-interpolation comparison tests selected spectral consequences of these
operations. It is not by itself an operator equality in an undeclared ambient
space.

\subsection*{Tight-binding model class}

Let
\begin{equation}
    \mathbf H_{\mathrm{TB}}(\mathbf k;\boldsymbol\theta)
    \in
    \mathbb C^{M\times M}
\end{equation}
be a member of a declared orthogonal $sp^3s^{\ast}$ model class. The parameter
domain fixes the orbital ordering, site convention, neighbor shells,
Slater--Koster structure, crystal symmetries, Fourier sign, spin convention,
and energy reference. Changing the permitted shell range or symmetry-allowed
terms changes the model class.

A tight-binding model is not defined solely by a parameter vector. Its
operator is
\begin{equation}
    \mathbf H_{\mathrm{TB}}(\mathbf k;\boldsymbol\theta)
    =
    \sum_{\mathbf R}
    e^{i\mathbf k\cdot\mathbf R}
    \mathbf H_{\mathrm{TB}}(\mathbf R;\boldsymbol\theta),
\end{equation}
where the retained blocks and Fourier convention are part of the model. A
nearest-neighbor model and a second-neighbor model belong to different
restricted classes even when their matrices have the same dimension at each
$\mathbf k$.

\subsection*{Direct KS-DFT-to-TB route}

The direct route determines a parameter vector
$\boldsymbol\theta_{\mathrm{dir}}$ from selected parent spectral data. An
abstract training objective is
\begin{equation}
    \mathcal L_{\mathrm{dir}}(\boldsymbol\theta)
    =
    \sum_{q\in\mathcal T_E}
    w_q\,
    \ell_q
    \left(
        \mathbf H_{\mathrm{TB}}(\boldsymbol\theta),
        \hat H_{\mathrm{KS}}^{(P)}
    \right),
\end{equation}
where $\mathcal T_E$ is the frozen spectral training set and every loss term
specifies its units and normalization.

This route can compare eigenvalues and derived band quantities without first
constructing a common orbital basis. It does not automatically supply an
operator Frobenius norm between the KS and TB Hamiltonians. Such a norm would
require an additional isometric embedding or alignment between the TB orbital
space and the retained KS subspace. Equal numbers of bands do not provide that
map.

The direct route must therefore report spectral fitting error separately from
withheld spectral and band-edge errors. A small training objective does not
establish agreement of orbital character, projectors, gauges, or matrix
elements.

\subsection*{Wannier-mediated route}

The mediated route first constructs a candidate Wannier operator and then
aligns the TB orbital coordinates with the Wannier coordinates. Let
\begin{equation}
    \mathbf C:
    \mathbb C^M_{\mathrm{TB}}
    \longrightarrow
    \mathbb C^M_{\mathrm W}
\end{equation}
be a symmetry-compatible unitary alignment after geometry, orbital identity,
spin, units, and energy reference have been checked. A proposed real-space
operator objective is
\begin{equation}
    \mathcal L_{\mathrm{med}}
    (\mathbf C,\boldsymbol\theta)
    =
    \sum_{\mathbf R}
    \omega_{\mathbf R}
    \left\lVert
        \mathbf H_{\mathrm W}(\mathbf R)
        -
        \mathbf C
        \mathbf H_{\mathrm{TB}}(\mathbf R;\boldsymbol\theta)
        \mathbf C^\dagger
    \right\rVert_{\mathrm F}^2.
\end{equation}
The direction of $\mathbf C$, the included blocks, shell weights,
normalization, and treatment of omitted blocks must be fixed before the
objective is interpreted.

The alignment does not repair incompatible subspaces. Principal angles,
overlap singular values, orbital correspondence, and alignment conditioning
must be checked before the operator residual is minimized. Optimizing
$\mathbf C$ over an unrestricted unitary group could hide physically
inadmissible orbital mixing; the admissible alignment class must preserve the
structures required by the $sp^3s^{\ast}$ interpretation.

\subsection*{Reciprocal-space and real-space Frobenius residuals}

For a common aligned gauge, define
\begin{equation}
    \Delta\mathbf H(\mathbf k)
    =
    \mathbf H_{\mathrm W}(\mathbf k)
    -
    \mathbf C
    \mathbf H_{\mathrm{TB}}(\mathbf k;\boldsymbol\theta)
    \mathbf C^\dagger
\end{equation}
and
\begin{equation}
    \Delta\mathbf H(\mathbf R)
    =
    \mathbf H_{\mathrm W}(\mathbf R)
    -
    \mathbf C
    \mathbf H_{\mathrm{TB}}(\mathbf R;\boldsymbol\theta)
    \mathbf C^\dagger.
\end{equation}
For a complete discrete Fourier pair with the convention above, Parseval's
identity gives
\begin{equation}
    \sum_{\mathbf k}
    \left\lVert
        \Delta\mathbf H(\mathbf k)
    \right\rVert_{\mathrm F}^2
    =
    N_{\mathbf k}
    \sum_{\mathbf R}
    \left\lVert
        \Delta\mathbf H(\mathbf R)
    \right\rVert_{\mathrm F}^2.
\end{equation}
This equivalence does not survive unchanged when nonuniform weights, incomplete
$\mathbf R$ sets, different reciprocal grids, or shell-dependent objectives
are introduced. Those choices define different losses and must be reported as
such.

The Frobenius norm is invariant under a common unitary coordinate change:
\begin{equation}
    \left\lVert
        \mathbf U^\dagger
        \Delta\mathbf H
        \mathbf U
    \right\rVert_{\mathrm F}
    =
    \left\lVert
        \Delta\mathbf H
    \right\rVert_{\mathrm F}.
\end{equation}
It is not invariant under independently changing only one operand. Gauge
invariance therefore follows only after the two representations and their
alignment map are transformed consistently.

\subsection*{The three pairwise relations}

The proposed bulk experiment contains three pairwise relations, but they answer
different questions:
\begin{description}
    \item[KS-DFT to Wannier.] Does the selected Wannier subspace and
      interpolation reproduce the declared retained parent information over
      the training and withheld domains? Relevant diagnostics include
      principal angles where a common ambient representation exists,
      interpolation errors, band-edge errors, centers, spreads, and real-space
      block decay.
    \item[KS-DFT to TB.] Does the direct TB model reproduce the frozen parent
      spectral and band-edge targets? Without an additional orbital embedding,
      this is primarily a spectral comparison rather than an operator
      Frobenius comparison.
    \item[Wannier to TB.] Can the declared TB model class approximate the
      aligned Wannier operator? This is the natural location for blockwise and
      shell-resolved Frobenius residuals.
\end{description}
These relations should not be collapsed into a single score. Their errors have
different meanings, domains, and dependencies.

\subsection*{Path dependence}

Let $\boldsymbol\theta_{\mathrm{dir}}$ be the direct-route result and
$\boldsymbol\theta_{\mathrm{med}}$ the mediated-route result. If both final TB
operators can be placed in one canonical orbital convention, define
\begin{equation}
    \Delta\mathbf H_{\mathrm{path}}(\mathbf k)
    =
    \mathbf H_{\mathrm{TB}}
    (\mathbf k;\boldsymbol\theta_{\mathrm{dir}})
    -
    \mathbf C_{\mathrm{path}}
    \mathbf H_{\mathrm{TB}}
    (\mathbf k;\boldsymbol\theta_{\mathrm{med}})
    \mathbf C_{\mathrm{path}}^\dagger,
\end{equation}
where $\mathbf C_{\mathrm{path}}$ records any required admissible alignment.
A proposed route discrepancy is
\begin{equation}
    D_{\mathrm{path}}^2
    =
    \sum_{\mathbf k\in\mathcal K_{\mathrm{cmp}}}
    w_{\mathbf k}
    \left\lVert
        \Delta\mathbf H_{\mathrm{path}}(\mathbf k)
    \right\rVert_{\mathrm F}^2.
\end{equation}

A parameter-vector difference is meaningful only if both vectors use the same
identifiable parameterization and canonical convention. Distinct parameter
vectors can represent the same operator, while similar parameter vectors can
produce different operators when conventions differ. The operator discrepancy
is therefore primary.

A nonzero path discrepancy is not automatically an error. It may reflect the
difference between a spectral objective and an operator objective, sensitivity
to the retained Wannier subspace, alignment constraints, model-class
limitations, optimization nonuniqueness, or numerical error. These causes must
be investigated separately.

\subsection*{Error accounting}

The comparison should retain the following error classes separately:
\begin{enumerate}
    \item \textbf{Parent-model error:} limitations of the declared Kohn--Sham
          physical approximation relative to independent physical evidence.
    \item \textbf{Parent numerical error:} incomplete convergence of the
          plane-wave, reciprocal-space, structural, or solver representation.
    \item \textbf{Subspace error:} sensitivity to retained rank, projections,
          and disentanglement windows.
    \item \textbf{Wannier representation error:} interpolation and finite-grid
          effects for the declared retained subspace.
    \item \textbf{Alignment sensitivity:} dependence on orbital matching,
          gauge, energy zero, and the admissible map.
    \item \textbf{TB model-reduction error:} inability of the prescribed
          orbital, symmetry, and neighbor-range class to reproduce the aligned
          parent representation.
    \item \textbf{Optimization and identifiability error:} uncertainty caused
          by incomplete minimization or nonunique parameter descriptions.
\end{enumerate}
No scalar total error is defined by this note. Combining these components would
require an owning rule that establishes compatible units, norms, correlations,
and validation domains.

\section*{Part II: Proposed numerical experiment}

\subsection*{Common parent and frozen comparison data}

Both routes must consume the same accepted pristine-silicon Kohn--Sham parent,
including the same physical branch, lattice, pseudopotential identity,
converged numerical settings, reciprocal conventions, and energy reference.
The parent dataset must be divided before fitting into frozen training and
withheld comparison sets.

The bulk pilot targets the valence and low-conduction information needed for
the orthogonal $sp^3s^{\ast}$ hierarchy, indirect Kohn--Sham gap, conduction
valley, and electron effective masses. The retained rank, windows, projection
choices, and interpolation domain must be supplied by the applicable accepted
records rather than selected retrospectively to improve the comparison.

No production parent is asserted to exist by this note. The following stages
are a proposed experiment to be performed only after the required parent and
its provenance are available.

\subsection*{Stage A: establish the retained parent target}

The first stage should record:
\begin{enumerate}
    \item the sampled $\mathbf k$ points and weights;
    \item the retained rank and target energy region;
    \item the training and withheld band data;
    \item the band-edge quantities and their extraction rules;
    \item the parent numerical uncertainty relevant to those quantities; and
    \item the representation in which subspace overlaps can be evaluated.
\end{enumerate}
This stage supplies the common target for both reduction routes. It does not
fit a TB model or select a Wannier gauge.

\subsection*{Stage B: construct and diagnose the Wannier representation}

A candidate Wannier90-derived representation should retain its initial
projections, inner and outer windows when applicable, localization settings,
centers, spreads, convergence history, Fourier conventions, and software
identity. Its evaluation should include:
\begin{enumerate}
    \item agreement with retained parent bands on the construction grid;
    \item interpolation errors at withheld $\mathbf k$ points;
    \item errors in the indirect gap, valley position, and effective masses;
    \item principal-angle or projector diagnostics where the common ambient
          data permit them; and
    \item decay of $\lVert\mathbf H_{\mathrm W}(\mathbf R)\rVert_{\mathrm F}$
          with lattice vector and neighbor shell.
\end{enumerate}
A candidate that fails the declared retained-subspace or interpolation criteria
should not be promoted to the mediated TB comparison merely because its
Wannier functions appear localized.

\subsection*{Stage C: direct spectral TB reduction}

The direct route should fit the first declared $sp^3s^{\ast}$ model class to
the frozen spectral training data. The objective, weights, parameter bounds,
symmetry constraints, initialization strategy, and optimization stability
checks must be fixed and retained.

After fitting, the model should be evaluated on the withheld band points and
on the declared indirect gap, valley position, and longitudinal and transverse
effective masses. Training and withheld errors must be shown separately. If the
nearest-neighbor class fails a prespecified criterion, the experiment may
proceed through the already declared hierarchy; the class sequence must not be
expanded after inspecting results merely to force acceptance.

\subsection*{Stage D: Wannier-mediated operator reduction}

The mediated route should align the declared TB orbital space with the accepted
Wannier space and fit or project the same ordered $sp^3s^{\ast}$ hierarchy
against the retained Wannier blocks. The experiment should record:
\begin{enumerate}
    \item the direction and admissible class of the alignment map;
    \item overlap, principal-angle, and conditioning diagnostics;
    \item the scalar energy-reference convention;
    \item included and omitted $\mathbf R$ blocks;
    \item shell weights and normalization;
    \item onsite, hopping, orbital-block, and symmetry-channel residuals; and
    \item sensitivity to the permitted neighbor range.
\end{enumerate}
The resulting model must also be evaluated on the same withheld spectral and
band-edge quantities used for the direct route. Operator fitting does not waive
spectral evaluation, and spectral agreement does not waive the operator
residual.

\subsection*{Stage E: compare the routes}

The final comparison should place the direct and mediated TB operators in one
canonical orbital, lattice, spin, Fourier, unit, and energy-reference
convention. It should then report:
\begin{enumerate}
    \item the route discrepancy $D_{\mathrm{path}}$ over a frozen reciprocal
          comparison set;
    \item corresponding real-space block differences where the complete
          Fourier conventions permit them;
    \item differences in withheld bands and band-edge quantities;
    \item model-class and neighbor-shell dependence;
    \item alignment and optimization sensitivity; and
    \item whether one parameter set satisfies both the prespecified spectral
          and operator criteria.
\end{enumerate}
If no tested model satisfies the declared criteria, the result should be
reported as failure of the tested reduction class rather than as failure of the
Kohn--Sham or Wannier parent.

\subsection*{Proposed comparison tables}

A first table should identify the immutable inputs and conventions shared by
both routes: parent identity, lattice, reciprocal coordinates, retained rank,
training set, withheld set, spin convention, units, and energy reference.

A second table should report the KS-to-Wannier diagnostics: interpolation
errors, band-edge discrepancies, projector or principal-angle information,
centers, spreads, and block-decay summaries.

A third table should compare direct and mediated TB models using separate
columns for training spectral error, withheld spectral error, indirect gap,
valley position, effective masses, aligned operator residual, model class, and
neighbor range. Blank or inapplicable cells should remain explicit rather than
being filled by incomparable quantities.

\section*{Plot placeholders}

\begin{enumerate}
    \item \textbf{Placeholder---route diagram.} Show the direct
          KS-DFT-to-TB route and the KS-DFT-to-Wannier-to-TB route, labeling
          subspace selection, gauge choice, alignment, spectral fitting, and
          operator reduction as distinct operations.
    \item \textbf{Placeholder---KS-to-Wannier interpolation.} Overlay retained
          parent and Wannier-interpolated bands on the frozen training and
          withheld reciprocal-space domains.
    \item \textbf{Placeholder---Wannier interpolation error.} Plot bandwise or
          subspace-aware interpolation discrepancies against $\mathbf k$,
          marking the indirect valley and withheld points.
    \item \textbf{Placeholder---real-space block decay.} Plot
          $\lVert\mathbf H_{\mathrm W}(\mathbf R)\rVert_{\mathrm F}$ against
          distance and neighbor shell, with onsite and symmetry-distinct blocks
          identified.
    \item \textbf{Placeholder---direct-route spectral comparison.} Overlay the
          parent and direct-fit TB bands, distinguishing training from withheld
          points.
    \item \textbf{Placeholder---mediated-route spectral comparison.} Overlay
          the parent, Wannier-interpolated, and operator-fit TB bands on the
          same reciprocal-space domain.
    \item \textbf{Placeholder---shell-resolved operator residual.} Plot
          $\lVert\Delta\mathbf H(\mathbf R)\rVert_{\mathrm F}$ by distance,
          shell, and selected orbital block.
    \item \textbf{Placeholder---model-class progression.} Plot spectral and
          operator discrepancies separately for the prespecified
          nearest-neighbor and subsequent tested model classes.
    \item \textbf{Placeholder---route discrepancy.} Plot
          $\lVert\Delta\mathbf H_{\mathrm{path}}(\mathbf k)\rVert_{\mathrm F}$
          over the frozen reciprocal comparison path or mesh.
    \item \textbf{Placeholder---band-edge comparison.} Display the indirect
          gap, valley position, and longitudinal and transverse effective-mass
          discrepancies for the Wannier, direct-TB, and mediated-TB
          representations without combining them into one score.
    \item \textbf{Placeholder---alignment sensitivity.} Show the change in the
          operator objective under the declared admissible alignment variants,
          together with principal-angle and conditioning diagnostics.
    \item \textbf{Placeholder---residual decomposition.} Display onsite,
          hopping, orbital-block, symmetry-channel, and omitted-range
          contributions to the mediated operator discrepancy.
\end{enumerate}
Each completed plot must identify the parent dataset, retained subspace,
training or withheld status, units, normalization, alignment convention, model
class, and evidentiary status. A placeholder is proposed work, not a calculated
result.

\section*{Interpretation boundary}

The proposed experiment asks whether a declared $sp^3s^{\ast}$ hierarchy can
preserve the selected bulk-silicon Kohn--Sham information through two different
reduction routes. The KS-to-Wannier relation tests retained representation and
interpolation. The direct KS-to-TB relation tests a constrained spectral fit.
The Wannier-to-TB relation tests an aligned operator reduction. The comparison
of their final TB operators tests path dependence.

Agreement between the routes would not establish that the Kohn--Sham parent is
an exact description of experimental silicon. Disagreement would not by itself
identify a single cause. Scientific interpretation requires the parent,
subspace, representation, alignment, model-class, optimization, and withheld
errors to remain visible as separate quantities.
